An increase in an incorrect-step score does not by itself imply a worse ranking. Ranking depends on which target receives the higher score, while the gap measures their separation. Wording can change the scores or narrow the gap without reversing their order. The preserved order in Figure~\ref{fig:gaps} shows that the model still favors the correct target within each constructed pair under the tested templates. This is a within-question comparison, not accuracy at a fixed classification threshold.

The changes relative to B cannot be attributed entirely to confidence. N also changes scores, although it makes no explicit claim of certainty. Compared with B, both templates add words and lowercase the original first letter of the target sentence. These post hoc contrasts combine wording content, input length, and capitalization changes; the experiment does not estimate their separate contributions.

The largest positive and largest negative incorrect-step effects in Figure~\ref{fig:individual} illustrate opposite responses to the same phrase replacement. These extreme cases are not representative of the sample. For the incorrect equality $3 \times 834 = 2503$, the score changes from 0.3047 under N to 0.4942 under A. For $182 - 124 = 59$, it changes from 0.5996 to 0.4210. Both targets use the same error rule, but this comparison does not identify why the wording affects them differently.

The design has several limits. It evaluates one model and one wording pair. Equal token counts do not isolate confidence from the meaning or familiarity of the specific words. The questions favor short arithmetic calculations, and every error is an upward change of one. More difficult reasoning errors, alternative confidence phrases, and other models may behave differently. The sample is also small and purposively selected. Three arithmetic expressions recur across different questions. Question-level resampling does not account for possible dependence between similar constructions.
